\section{New theorem: periodic Alexandrov existence for every Radon target}
All the real arguments in this section are proved here. They do not assume
a periodic transport or Monge--Ampere existence theorem. Write
$T=\mathbb R^2/\mathbb Z^2$, $F=[0,1)^2$, $\rho(r)=|r|^2/2$, and let
$m$ be Lebesgue probability measure on $T$. For a finite convex function
$\Phi=\rho+\bar q$ with continuous periodic $\bar q$, its conjugate is
$\Psi(y)=\sup_r(r\cdot y-\Phi(r))$.

\begin{lemma}\label{newrad:convex}
Both conjugates are finite, mutually conjugate, and satisfy
\[
 \Phi(r+k)=\Phi(r)+r\cdot k+|k|^2/2,\qquad
 \Psi(y+k)=\Psi(y)+y\cdot k+|k|^2/2.
\]
Every subgradient $p\in\partial\Phi(r)$ satisfies
$|p_i-r_i|\le1/2$. Consequently $\bar q$ is $1/2$-Lipschitz in each
coordinate; if $\max_T\bar q=0$, then $-1/2\le\bar q\le0$.
Convex functions here are differentiable Lebesgue almost everywhere.
\end{lemma}
\begin{proof}
A bounded periodic perturbation of $\rho$ is superlinear, so conjugate
suprema are attained on compact sets. Substitution $y=z+k$ in the
conjugate proves its displayed translation law. A supporting plane to
the epigraph at $(r,\Phi(r))$ gives a nonempty subgradient: separation
of the epigraph from $(r,\Phi(r)-\epsilon)$ and a convergent normalized
separating-plane subsequence prove this assertion; the vertical coefficient
is nonzero since $\Phi$ is bounded on a neighborhood of $r$.
The subgradient inequality at $r\pm e_i$ gives
$r_i-1/2\le p_i\le r_i+1/2$. More generally a finite convex function
has bounded subgradients on each compact interior set by applying its
supporting inequalities at a fixed surrounding ball. These inequalities
also prove local Lipschitz continuity.

On every coordinate line, left and right derivatives of a convex
function are monotone and disagree at only countably many points:
each strict gap contains a distinct rational number. Fubini therefore
gives existence of both coordinate derivatives almost everywhere.
There every supporting vector has these two coordinates, so the
subgradient is a singleton. This implies differentiability: any
subgradient at a sequence tending to $r$ has bounded subsequences;
passing its supporting inequalities to the limit puts each cluster
in $\partial\Phi(r)$, and the two supporting inequalities at the
endpoints of an increment sandwich the increment by their slopes.
Along a coordinate line, $\bar q$ is absolutely continuous on compact
intervals, and its almost-everywhere derivative is between $-1/2$ and
$1/2$. Integration proves the Lipschitz bound. Two points of $T$ have
coordinate distances at most $1/2$, giving oscillation at most $1/2$.

Finally $y\in\partial\Phi(r)$ if and only if
$\Phi(r)+\Psi(y)=r\cdot y$, if and only if
$r\in\partial\Psi(y)$. The supporting plane just obtained gives
$\Psi^*(r)=\Phi(r)$, proving mutual conjugacy without a biconjugation
theorem assumption.
\end{proof}

\begin{lemma}\label{newrad:Alex}
The periodic real Alexandrov measure of $\Phi$ is the probability
\[
 M_\Phi=(\nabla\Psi\bmod\mathbb Z^2)_*m.
\]
It agrees with Lebesgue measure of subgradient images in the ordinary
Alexandrov convention, without a factorial. It is continuous under
uniform convergence of periodic $\bar q$. For smooth $\Phi$ it equals
$\det(\operatorname{Hess}\Phi)\,dr_1dr_2$ on $T$.
\end{lemma}
\begin{proof}
The gradient is measurable on its full-measure differentiability set
(use limits of difference quotients); arbitrary values on its null
complement do not affect its pushforward. Translation equivariance
makes the expression well-defined. For a Borel $E\subset T$, choose
its unique representatives $\widetilde E\subset F$. Outside the
nondifferentiability set of $\Psi$ and all its lattice translates,
\[
 \sum_{k\in\mathbb Z^2}
 1_{\partial\Phi(\widetilde E)}(y+k)
 =1_E(\nabla\Psi(y)\bmod\mathbb Z^2).
\]
Indeed the unique point related to $y+k$ is $\nabla\Psi(y)+k$, and
exactly one translate belongs to $F$. The supporting-plane equivalence
in the preceding proof identifies this relation with the subgradient
image. The image differs from the measurable inverse-gradient set
only by a subset of a Lebesgue null set, so is Lebesgue measurable in
the completed measure. Integrating and summing over translates proves
the Alexandrov formula, including at atoms on the cut of $F$. This
also proves countable additivity rather than assuming it for images.

If $\bar q_n\to\bar q$ uniformly, then
$|\Psi_n-\Psi|\le\|\bar q_n-\bar q\|_\infty$ globally. Local
boundedness of subgradients follows from the preceding surrounding-ball
argument uniformly in $n$. At a differentiability point of $\Psi$,
any cluster of $\nabla\Psi_n(y)$ belongs to $\partial\Psi(y)$ by
passing supporting inequalities to the limit. Thus these gradients
converge wherever all of them and $\Psi$ are differentiable, a common
full-measure set. Dominated convergence applied to any continuous
periodic test proves $M_{\Phi_n}\rightharpoonup M_\Phi$.

For smooth positive-definite Hessian, replace $\Phi$ if necessary by
$(1-\epsilon)\Phi+\epsilon\rho$, whose Hessian is at least
$\epsilon I$. Its gradient is a global diffeomorphism: strong
convexity gives injectivity, minimizing $\Phi(r)-y\cdot r$ gives
surjectivity, and the inverse function theorem gives a smooth inverse.
It commutes with lattice translations, so induces a torus
diffeomorphism. Change of variables for this inverse proves the
determinant formula. For a smooth merely nonnegative Hessian use the
same strict perturbations; the continuity just proved and uniform
convergence of their determinants prove the formula by taking
$\epsilon\downarrow0$.
\end{proof}

\begin{theorem}\label{newrad:realexist}
For every Radon probability $P$ on $T$ there is a continuous periodic
$\bar q$, unique up to a constant, with $\Phi=\rho+\bar q$ convex
and $M_\Phi=P$. The normalized solution has range in $[-1/2,0]$.
No absolute continuity, density bound, or absence of atoms is required.
\end{theorem}
\begin{proof}
First let $P=\sum_{i=1}^N m_i\delta_{x_i}$, with distinct torus
points, $m_i>0$, and $\sum_i m_i=1$. Fix lifts $x_i\in F$ and put
\[
 c_i(y)=\tfrac12\min_{k\in\mathbb Z^2}|y-x_i-k|^2,
 \qquad H(w)=\int_T\max_i(w_i-c_i(y))\,dm(y)-\sum_i m_iw_i.
\]
The $c_i$ are continuous periodic functions between $0$ and $1/4$.
$H$ is continuous convex and invariant under addition of a common
constant to all $w_i$. On the gauge $\max_i w_i=0$,
\[
 H(w)\ge-1/4+\sum_i m_i(-w_i).
\]
Its sublevel $H\le H(0)\le0$ is compact, since every $w_i$ lies
between $-1/(4m_i)$ and $0$. It contains $0$, so $H$ has a global
minimum $w$ there; shifting any vector to this gauge shows it is a
minimum in all $\mathbb R^N$.

On a compact set in the $y$-plane, only finitely many nearest lattice
copies of each $x_i$ are relevant. Equal weighted squared-distance
functions for distinct copies have affine hyperplane equality sets.
Their union over all copies is countable and Lebesgue null. Thus
almost every $y$ has a unique winning site and copy. Difference
quotients of the maximum are bounded by $1$ in each weight coordinate;
dominated convergence gives
\[
 \partial_{w_i}H= m(C_i)-m_i,
 \quad C_i=\{y\in T:i\text{ is the winning site}\}.
\]
At the minimum all derivatives vanish. Hence every cell has its
prescribed positive volume; no cell positivity is presumed beforehand.

Define
\[
 \Psi_w(y)=\rho(y)+\max_i(w_i-c_i(y))
 =\max_{i,k}\big((x_i+k)\cdot y-|x_i+k|^2/2+w_i\big).
\]
This is finite convex, locally a maximum of finitely many planes, and
quasi-periodic with the translation law of Lemma~\ref{newrad:convex}.
Its conjugate $\Phi_w$ has the same law, and is $\rho$ plus a bounded
periodic continuous function. At almost every point of $C_i$ its
gradient is a lift $x_i+k$. Therefore
$M_{\Phi_w}=(\nabla\Psi_w\bmod\mathbb Z^2)_*m=P$.
This uses the conjugate in the required direction: it is $\Phi_w$,
not $\Psi_w$, whose Alexandrov measure is the target.
For additional cell precision, the full cell of $x_i+k$ is the
intersection of the plane-winning halfspaces. It is a closed convex
polytope contained in $x_i+k+[-1/2,1/2]^2$, since this copy must beat
its own neighboring translates. Cells of successive lattice copies
are translates. Summing their intersections with $F$ proves that the
volume of each full cell equals $m_i$. Its boundary is null by the
hyperplane argument. Fenchel equality identifies this full cell with
$\partial\Phi_w(x_i+k)$.

For general $P$, partition $T$ into a finite grid of half-open squares
of diameter at most $\sqrt2/n$; place each square's mass at one point
in it and omit zero masses. This gives finite probabilities $P_n$ and
$|\int h\,d(P_n-P)|\le\omega_h(\sqrt2/n)$ for every continuous $h$.
Take the atomic solutions and normalize their maxima to zero. The
uniform range and coordinate Lipschitz bounds of Lemma~\ref{newrad:convex}
give a uniformly convergent subsequence. Explicitly, successive finite
grids and bounded scalar subsequences give a diagonal subsequence
Cauchy on a dense union of grids; the Lipschitz bound upgrades it to
uniform Cauchyness. Its continuous periodic limit $\bar q$ has convex
$\rho+\bar q$. Lemma~\ref{newrad:Alex} gives $M_\Phi=P$.

Here is uniqueness, proved without a transport uniqueness theorem.
Let $(\Phi_j,\Psi_j)=(\rho+q_j,\rho+h_j)$ be two solutions and
$T_j(y)=\nabla\Psi_j(y)$, with $T_j\bmod\mathbb Z^2$ pushing $m$
to $P$. Fenchel's inequality at all lattice lifts gives
\[
 q_j(x)+h_j(y)+c(x,y)\ge0,
 \quad c(x,y)=\tfrac12\min_k|x-y-k|^2.
\]
By the subgradient bound, $T_j(y)-y$ has both coordinates in
$[-1/2,1/2]$, so its squared Euclidean cost is exactly this torus
cost. Fenchel equality holds for $(T_j(y),y)$. Thus the cost of the
plan $(y,T_j(y)\bmod\mathbb Z^2)_*m$ is
$-\int q_j\,dP-\int h_j\,dm$. Applying the other feasible
inequality to each plan shows that these two costs are equal.
The nonnegative slack for the second potentials on the first plan
therefore vanishes almost everywhere. Using the actual nearest lift
$T_0(y)$, this means $T_0(y)\in\partial\Psi_1(y)$.
At the common differentiability points it forces
$\nabla\Psi_0(y)=\nabla\Psi_1(y)$ as real vectors, not just modulo
lattice translations. The Lipschitz difference $\Psi_0-\Psi_1$ has
zero weak gradient: this follows by absolute continuity on coordinate
lines and Fubini. Convolution makes its gradient zero pointwise, so
each convolution is constant on the connected plane. Local uniform
convergence makes the difference constant. Conjugation now proves
that $q_0-q_1$ is constant.
\end{proof}

\section{Actual analytic lift, the entire full dual, and bounded energy}
This section uses precisely the conventional CLD/GR foundation boundary
of the fully retained 421 reconstruction, with its fifteen immediate
statement-and-proof units independently reread here. It does not certify
genuine nonalgebraization or coherent descent. The given analytic quotient
is used as input. No primary theorem on toric model lattices or
algebraization, and no Vilsmeier comparison theorem, is used.

\begin{theorem}\label{newrad:analytic}
On the designated $K=\mathbb C((s))$ quotient $X$, for every positive
Radon measure $\mu$ of raw mass $2$ supported on
$S=\sigma(T)$, there exists a normalized continuous invariant weight
$q=\bar q\circ\tau$ in the \emph{original entire} class $\C$ with
\[
 \MA(q)=\mu,\quad -1/2\le q\le0,\quad -1\le E(q)\le0,
 \quad D(\mu)=E(q)-\int q\,d\mu\in[0,1].
\]
Thus the full dual is attained, and $q$ has an actual continuous
decreasing presentation in the original $\mathcal G$.
The result is at the disclosed conventional foundation boundary.
\end{theorem}
\begin{proof}
Since $S$ is the continuous image of compact $T$ into Hausdorff $X$,
it is closed and $\sigma:T\to S$ is a homeomorphism, inverse to
$\tau|_S$. Put $P=\tau_*\mu/2$. The preceding theorem supplies
$\bar q$ with $M_{\rho+\bar q}=P$.

Choose even nonnegative smooth approximate identities on $\mathbb R^2$
and convolve the periodic $\bar q$. Set $\bar q_j=\bar q*\chi_j$.
Uniform continuity gives global uniform convergence on $T$.
Because the kernel has mean zero,
$\rho+\bar q_j=(\rho+\bar q)*\chi_j-\int\rho\chi_j$ is convex.
Hence $I+\operatorname{Hess}\bar q_j\ge0$.
On every injective quotient chart $\rho+\bar q_j\circ\tau$ is an
ordinary CLD smooth potential with positive Hessian: its tropicalization
uses the actual invertible coordinate functions $z_1,z_2$. Under a
deck shift $k$, the local potential changes by $r\cdot k+|k|^2/2$,
a smooth affine unit-logarithm combination with zero $\DD$.
Periodicity makes $q_j=\bar q_j\circ\tau$ a global continuous
function on the supplied quotient. The entire minimal continuous
sheaf of CLD8.1.7--8 and smooth-null-change rule8.8.5 therefore give
$q_j\in\C$ and, by uniform closure, $q\in\C$. This is an actual
uniform approximation by positive generators, not an assigned positivity
for a Borel limit.

For each $j$ the coordinate Hessian square is supported on $S$.
For clarity, at an off-Gauss point a Laurent polynomial has strict
cancellation below its monomial maximum. That strict inequality
persists on an ordinary neighborhood. All coordinate tropical values
there must have a tie among maximum monomials, hence lie in the finite
tropical hypersurface of that polynomial, of dimension at most one.
The coordinate top-form dimension rule makes the Hessian square zero
there. Equivalently this is the coordinate support rule in CLD8.8.6.
For the identity moment on an injective torus chart the tangent lattice
is saturated and the degree of the identity map is one. CLD14.5.2--3
calibrates its full open two-cells by the primitive Euclidean volume
vector. CLD5.10.1(c),(d) and5.10.2 give its top Radon pushforward and
reconstruction. Direct expansion of the square of the $2\times2$
coordinate Hessian gives the factor $2!$ (including both cross terms).
Consequently, as a measure on the actual analytic space,
\[
 \MA(q_j)=2\sigma_*\bigl[\det(I+\operatorname{Hess}\bar q_j)\,m\bigr]
 =2\sigma_*M_{\rho+\bar q_j}.
\]
The same reasoning for $\bar q=0$ gives $\theta^2=2\sigma_*m$.
The support argument is local in ordinary quotient charts; degree-one
calibration and periodic gluing introduce no additional field or
lattice multiplicity.

Real Alexandrov continuity proved above shows that the displayed
measures converge to $2\sigma_*P$ against \emph{every} continuous
analytic test $h$, since $h\circ\sigma\in C(T)$. Independently,
CLD8.5.1/8.5.5/8.7.1--2 imply that the actual recursive products
$\MA(q_j)$ converge as currents to the already defined $\MA(q)$.
Positive top currents are Radon by CLD7.2.8; the density argument in
its proof upgrades current convergence to all continuous compact tests.
Alternatively the explicit measures have common mass2, so approximation
of any continuous test by CLD smooth tests makes this upgrade directly.
Thus $\MA(q)=2\sigma_*P=\mu$. The operator is not defined by its
projection: the actual recursive operator and all analytic tests agree.

The real range bound is a bound on all $X$, because $q$ is invariant.
Energy monotonicity and $E(c)=2c$ give $-1\le E(q)\le0$.
Tangent concavity in the \emph{entire} class $\C$, reconstructed
by all-stage Green, gives
$E(a)\le E(q)+\int(a-q)\MA(q)$ for every $a\in\C$.
Taking the full supremum proves attainment at $q$.
For every normalized $a\in\C$, Gauss dominance and the Gauss range
bound give $\int(-a)\,d\mu\le1$, while $E(a)\le0$; hence
$D(\mu)\le1$. The zero competitor gives $D(\mu)\ge0$.
This also independently verifies finiteness of the full, not invariant,
dual for every skeletal Radon target.

For an explicit actual decreasing presentation, pass to convolution
indices with $\|q_j-q\|_\infty\le2^{-j}$ and put
$b_j=q_j+3\cdot2^{-j}$. The upper bound for $b_{j+1}$ and lower
bound for $b_j$ are respectively $q+2\cdot2^{-j}$ and the same
function. Thus $b_{j+1}\le b_j$, $b_j\downarrow q$ on all $X$,
and their energies have a finite lower bound. The constant presentation
of $q$ is another such presentation.
\end{proof}

\begin{corollary}\label{newrad:fast}
For every such skeletal target and every fast maximizing family in the
421 construction, the actual canonical weak measure is $\nu=\mu$,
$K=J=0$, $e=E(q)$, and $R=q$. Its Borel limit equals $q$ almost
everywhere for $\mu$, every bounded MA measure and every bounded mixed
measure, and at every full point. This assertion does not assign a
new psh class to that Borel limit at its remaining points.
\end{corollary}
\begin{proof}
Full attainment gives $I(a_n,q)\le4\varepsilon_n$ by the exact
midpoint formula. The retained energy and continuous-test moduli give
$e=E(q)$ and $\nu=\MA(q)=\mu$; full-point atoms give equality of
full-point limits. The all-affinoid Shilov formula with these values
and continuity of $q$ gives $R=q$. Since $u\le R$ and
$\int u\,d\mu\ge e-D(\mu)=\int q\,d\mu$, equality and $J=0$
follow. The fixed bounded-test moduli and fast tails give the other
almost-everywhere equalities. These are now consequences of proved
existence, rather than unproved attainment assumptions.
\end{proof}

\section{Uniqueness against every actual generated presentation}
\begin{theorem}\label{newrad:Gunique}
At the same explicit bounded/generated foundation boundary, a bounded
solution $q\in\C$ is unique modulo constants among \emph{all}
actual weights $v\in\mathcal G$ having $\MA(v)=\MA(q)$.
In particular the skeletal theorem has full generated, not only
invariant, uniqueness. This does not construct a generated solution
for an arbitrary nonskeletal measure.
\end{theorem}
\begin{proof}
We give the generated calibration argument used here. Let
$c_n\downarrow v$ be an actual continuous presentation with
$E(c_n)$ bounded below, and subtract a common constant so $c_n\le0$.
Monotonicity makes its limiting energy finite. For two presentations
of the same limit, each fixed continuous term of one plus $\epsilon$
strictly exceeds the limit of the other. The open sets expressing
this strict inequality cover compact $X$, so some term of the other
is below that fixed term plus $\epsilon$ on all $X$. Monotonicity
and $E(a+\epsilon)=E(a)+2\epsilon$ prove presentation independence.
The same argument applies to every fixed finite convex combination
of presentations. Its energies are polynomials of degree at most3
in the simplex parameters. Values on a finite unisolvent set converge;
thus their coefficients converge. Such a set exists by elementary
linear algebra since a polynomial vanishing on an open simplex is zero.

Write $e=\lim E(c_n)$, $A_n(t)=E(tc_n)$ and
$P_n(t)=E((1-t)c_n+tw)$ for any fixed $w\in\C$.
First variation gives
$\int w\MA(c_n)=P'_n(0)+A'_n(1)$. Coefficient convergence and
the dense span of $\C$ give the canonical weak measure $\lambda$
of mass2. In the original generated operator this is $\MA(v)$.
Also $d=\lim\int c_n\MA(c_n)=\lim A'_n(1)$ is finite.
Passing the bounded tangent inequality for each fixed $w$ gives
$D(\lambda)\le e-d<\infty$. Applying its full dual inequality
to $c_n$ and using monotone convergence shows $v$ is
$\lambda$-integrable.

For $l\ge n$, the ordered cubic cocycle gives
\[
 0\le\int(c_n-c_l)\MA(c_l)
 \le3\bigl(E(c_n)-E(c_l)\bigr).
\]
With $n$ fixed, pass $l\to\infty$: the first term converges against
the fixed continuous $c_n$, and the diagonal term converges to $d$.
Then take $n\to\infty$, using monotone convergence and the proved
integrability. It follows that $d=\int v\,d\lambda$.
The presentation now gives the reverse inequality to the tangent bound,
so $D(\lambda)=e-\int v\,d\lambda$ and
$F_\lambda(c_n)\to D(\lambda)$.

The decreasing suprema $s_n=\sup_Xc_n$ have a finite lower bound:
$E(c_n)\le2s_n$. Their limit $s$ equals $\sup_Xv$. For example,
the nested nonempty compact sets $\{c_n\ge s\}$ have nonempty
intersection and provide a point where $v=s$; the reverse bound is
immediate. The normalized $a_n=c_n-s_n$ converge at every point to
$v-s$, have bounded lower energies, and still maximize $D(\lambda)$.
Normalize $q$ as well. Tangent concavity makes it a maximizer for
$\lambda=\MA(q)$. The midpoint estimate and full-point moduli give
$a_n(x)\to q(x)$ at every full graded-dimension-two point, whence
$v(x)-s=q(x)$ there.

The set of these full points is dense in the ordinary topology:
each nonempty ordinary open contains a pure two-dimensional affinoid
neighborhood with its nonempty finite Shilov boundary inside that
open, and GR6.3 says these points are full. The function $v-s$ is
USC as a decreasing limit of continuous functions. If at any point
$v-s<q$, USC and continuity would give a nonempty ordinary open of
strict inequality, contradicting the dense full-point equality.
Thus $v-s\ge q$ everywhere, including at possible poles.
For the converse, fix a local reference and an affinoid neighborhood
$A$ of any $x$. GR6.4 for each $\rho+a_n$ and the finiteness of
$\Gamma_A$ give
\[
 \rho(x)+v(x)-s\le
 \max_{\xi\in\Gamma_A}(\rho(\xi)+q(\xi)).
\]
Choose $A$ inside successively smaller ordinary neighborhoods and use
continuity of $\rho+q$. The right side can be made arbitrarily close
to $\rho(x)+q(x)$. Hence $v-s\le q$ everywhere. This proves actual
pointwise equality on all $X$ and uniqueness modulo constants.
No upper regularization, GR6.5 recognition, or historical all-curve
uniqueness claim is a premise of this argument.
\end{proof}

\section{New arbitrary-support consequences and the remaining existence gate}
\begin{theorem}\label{newrad:projection}
For every eligible arbitrary-support $\mu$, put
$\lambda=\sigma_*\tau_*\mu$. Then
\[
 D(\mu)\ge D(\lambda),\qquad
 D(\mu)=D(\lambda)\ \Longleftrightarrow\ \mu=\lambda.
\]
Consequently an invariant continuous weight attains the full dual of
$\mu$ if and only if $\mu$ is supported on $S$. This drops the
smoothness and strict-coordinate-Hessian assumptions of the earlier
attainment slice. For nonskeletal eligible $\mu$, there is an actual
$a\in\C$ and a positive certified gap as follows. Let $q_\lambda$
be the normalized skeletal solution, normalize $a$, and set
\[
 \Delta=\int a\,d(\lambda-\mu)>0,
 \quad C_0=\max(1,-E(a),-E(q_\lambda)).
\]
Then
\[
 D(\mu)-D(\lambda)\ge
 \min\{\Delta/2,\Delta^2/(96C_0)\}>0.
\]
This is a theorem on arbitrary-support input measures, not an equation
identifying those measures with their projection.
\end{theorem}
\begin{proof}
Gauss dominance gives $\int a\,d\mu\le\int a\,d\lambda$ for every
$a\in\C$. Thus $F_\mu(a)\ge F_\lambda(a)$ and the dual inequality
follows. Since $q_\lambda$ is invariant,
$F_\mu(q_\lambda)=F_\lambda(q_\lambda)=D(\lambda)$.
If the duals are equal it attains $D(\mu)$. The actual segment from
$q_\lambda$ to any $a\in\C$ has nonpositive right derivative at
this maximum. Its derivative is
$\int(a-q_\lambda)\,d(\lambda-\mu)$.
The invariant term has zero integral difference, so this gives the
reverse inequality to Gauss dominance. Integrals agree on $\C$,
and density of its real span gives $\mu=\lambda$.

Suppose any invariant $b\in\C$ attains $D(\mu)$. Its integrals
against $\mu$ and $\lambda$ agree, so
$D(\mu)=F_\lambda(b)\le D(\lambda)\le D(\mu)$.
The preceding equivalence forces skeletal support. Conversely the
skeletal existence theorem supplies an invariant attained weight.
If $\mu\ne\lambda$, density of the span gives some $a\in\C$ with
unequal integrals, and dominance makes their difference $\Delta>0$.
For the actual segment $q_t=(1-t)q_\lambda+ta$, all its normalized
endpoint energies are at least $-C_0$. The reconstructed Green bounds
$Q_{\theta_b}(a),Q_{\theta_b}(q_\lambda)\le6C_0$ for the endpoint
currents give $Q_{\theta_{q_t}}(a-q_\lambda)\le24C_0$.
Thus $|F_\mu''(q_t)|\le48C_0$. Twice integration gives
$F_\mu(q_t)\ge D(\lambda)+t\Delta-24C_0t^2$.
Use $t=\min(1,\Delta/(48C_0))$ to get the displayed lower bound.
The direction is an actual nonskeletal-sensitive CLD competitor, not
an arbitrary smooth perturbation declared positive.
\end{proof}

\begin{corollary}\label{newrad:nonskeletal-example}
There is a concrete eligible nonskeletal atomic input on the original
analytic carrier. Let $x$ be the Gauss point of the bidisc
$z_i=1+s u_i$, $|u_i|\le1$. Then $\tau(x)=0$, $x\ne\sigma(0)$,
and $\mu=2\delta_x$ has finite full dual. In fact
\[
 D(2\delta_x)\ge\frac16+\frac1{25165824}.
\]
It has no continuous invariant full-dual maximizer. Existence of a
noninvariant solution for this input is not proved here.
\end{corollary}
\begin{proof}
The coordinate initials of $u_1,u_2$ are independent, so $x$ has
full graded dimension2. Both original coordinate norms are1, while
$|z_i-1|(x)=|s|=e^{-1/8}$ differs from its norm1 at $\sigma(0)$.
The full-point atom construction gives $\MA(w_x)\ge\alpha\delta_x$.
Here its coordinate representative can be taken to be
$((z_1-1)/s,(z_2-1)/s)$, degree1 near the target Gauss point;
the fixed log-max fibre formula also proves this local atom directly.
To verify eligibility, put $B_x=D(\MA(w_x))<\infty$, given by the
bounded tangent formula. For normalized $a$, put
$A=\int(-a)\theta^2\le1$, $q_0=Q_\theta(a)$ and
$q_1=Q_{\theta_a}(a)$. Integrating
$E''(ta)=-2[(1-t)q_0+tq_1]$ twice gives
$C=-E(a)=A+2q_0/3+q_1/3$ and
$-E(ta)=tA+t^2(1-t/3)q_0+t^3q_1/3$.
Nonnegativity gives $q_0+q_1/3\le3(C-A)/2$; hence
\[
 -E(ta)\le t+\tfrac32t^2C,\qquad
 \int(-a)\MA(w_x)\le1+B_x+\sqrt{6B_x C}.
\]
For the second inequality apply the full dual to $ta$, divide by $t$,
and use $t=\min(1,\sqrt{2B_x/(3C)})$, with limiting conventions
when $B_x$ or $C$ is zero. Thus
$F_{2\delta_x}(a)\le-C+(2/\alpha)(1+B_x+\sqrt{6B_xC})$,
whose supremum over $C\ge0$ is finite. This proves finite FULL dual.

In the injective tropical ball $|r|<1/4$ around0, set
\[
 \psi=\log\max(1,|z_1-1|/e^{-1/16}),\qquad
 w=\max(0,1/256+\psi/64-\rho).
\]
On its closure $0\le\psi\le5/16$, so $w\le9/1024$.
On $|r|\ge3/16$ the reference satisfies $\rho\ge9/512$,
and hence $w=0$ on this outer collar. The actual log-max ramp,
maximum and zero-gluing argument puts $w$ in $\C$ globally.
At $\sigma(0)$ and $x$ its values are respectively $5/1024$ and
$4/1024$. Normalize $a=w-\sup_Xw$; its range is contained in
$[-9/1024,0]$, so $-E(a)\le9/512<1$.
With $\lambda=2\delta_{\sigma(0)}$ one has
$\Delta=\int a\,d(\lambda-\mu)=1/512$ and $C_0=1$.
The projection theorem and the exact atomic value $D(\lambda)=1/6$
give the claimed lower bound $\Delta^2/96=1/25165824$.
\end{proof}

\begin{proposition}\label{newrad:stability}
Normalized skeletal solutions depend continuously in the uniform norm
on weak convergence of their torus probabilities. If $W_1$ uses torus
coordinate $\ell^1$ distance and $\lambda_i=2\sigma_*P_i$, then
\[
 |D(\lambda_0)-D(\lambda_1)|\le W_1(P_0,P_1),\qquad
 I(q_0,q_1)\le8W_1(P_0,P_1).
\]
\end{proposition}
\begin{proof}
The normalized potentials are uniformly bounded and equi-Lipschitz.
Every uniformly convergent subsequence has, by the proved real MA
continuity, the limit target. Real uniqueness identifies the limit;
compactness therefore gives uniform convergence of the whole sequence.
Each $q_i$ is $1/2$-Lipschitz for coordinate $\ell^1$ distance.
For any coupling of $P_0,P_1$, its integral difference against the
raw measures is bounded by its transportation cost. Evaluating the
attained full objectives at $q_0,q_1$ proves the first inequality,
then taking the infimum over couplings defines $W_1$. The deficit of
$q_1$ for $\lambda_0$ is at most $2W_1$; the midpoint inequality
with $q_0$ gives the second bound. No existence of an optimal coupling
is needed for either estimate.
\end{proof}

\section{An exact atomic calibration and scoped verdict}
For the target $2\delta_{\sigma(0)}$, let
$\xi_i=r_i-\lfloor r_i\rfloor$ and
\[
 \bar q(r)=\tfrac12\xi_1(1-\xi_1)
             +\tfrac12\xi_2(1-\xi_2)-1/4.
\]
This is periodic continuous, normalized, and ranges in $[-1/4,0]$.
$\rho+\bar q$ is the sum of the piecewise-affine interpolants of
$r_i^2/2$ at integer $r_i$, minus $1/4$. Each one-dimensional
derivative jumps by1 at every integer, so its two-dimensional
Alexandrov measure is the unit lattice-point measure, descending to
$\delta_0$. The analytic theorem gives the actual equation.
The reference, mixed, and MA moments are respectively $-1/6,-1/3,-1/2$:
the reference is $2m$, the mixed measure is the sum of the two
coordinate-circle Lebesgue measures, and the MA measure is $2\delta_0$.
Hence $E(q)=-1/3$ and $D(2\delta_{\sigma(0)})=1/6$.
The mixed assertion can also be obtained by the same convolution:
the coefficient traces of the separable smooth Hessians converge to
the two circle measures, so it is an actual recursive mixed product.
Translations give the corresponding formula at every Gauss point.

Every-Gauss-Radon existence and full-dual attainment are new proved
sector deductions on the input analytic quotient at the disclosed
conventional source boundary. The real existence and continuity
arguments are complete and independent of an external Alexandrov or
optimal-transport theorem. Generated uniqueness against bounded
solutions has a new direct proof. The strict projection gap is a
new assertion for every eligible arbitrary-support measure.
The proposal does not remove the skeletal restriction from existence:
nonordered actual representation, legitimate arbitrary-support
variations/envelopes, and general full attainment remain unproved.
Genuine nonalgebraization, coherent descent, and the deeper recursive
source/book audit remain unproved here. The whole historical companion
is retained, but was not read or independently audited in its entirety.
No complete original P16.7 solution, publication or acceptance is
claimed. The original residual file is byte-unchanged and candidate
is false.
